%!TeX root=MemoriaTFG.tex

\chapter{Ablación: modelo fundacional inglés (DermLIP) con
traducción ES--EN}
\label{cap:anexoi}

Este anexo documenta una ablación independiente que cuantifica el
\emph{trade-off} entre la rama contrastiva de Dermapixel R0
defendida en la
sección~\ref{sec:spanderm-clip} y un \emph{wrapper} alternativo
basado en un modelo fundacional dermatológico inglés (DermLIP v2)
con traducción de etiquetas ES--EN, sobre el mismo conjunto de test
\emph{holdout} de DermapixelAI~1.0.

\section{Motivación}
\label{sec:anexoi-motivacion}

Una pregunta legítima planteada durante la revisión paper-grade
de la rama contrastiva de Dermapixel R0 es si el esfuerzo experimental
documentado en la sección~\ref{sec:spanderm-clip} (cinco
iteraciones de entrenamiento cross-modal castellano) se justifica
frente a una alternativa pragmática: emplear un modelo fundacional
dermatológico inglés ya alineado imagen-texto, sin entrenar ningún
modelo adicional, sobre las etiquetas castellanas traducidas al
inglés. Esta ablación cuantifica el \emph{trade-off} entre ambas
estrategias sobre el mismo conjunto de test \emph{holdout} de
DermapixelAI~1.0.

\section{Protocolo experimental}
\label{sec:anexoi-protocolo}

El modelo fundacional inglés evaluado es DermLIP v2
(\texttt{PanDerm-base-w-PubMed-256}), un sistema CLIP-style con
codificador visual PanDerm-Base y codificador textual PubMedBERT,
preentrenado contrastivamente sobre el dataset
Derm1M~\cite{derm1m} ($\approx 403\,000$ pares imagen-texto).
Ambas torres permanecen congeladas; no se entrena ningún
parámetro. Las $54$ clases L3 representadas en el test se
tradujeron al inglés clínico mediante un glosario estándar
(Bolognia/Rook) con verificación cruzada de unicidad, y por cada
clase se generó un \emph{rich prompt} clínico en inglés ($\sim 37$
palabras, longitud equiparable a los \emph{prompts} castellanos
de Dermapixel R0). La inferencia replica el protocolo híbrido de
la sección~\ref{sec:spanderm-clip} (prototipos visuales sobre
\emph{train} más texto rico, régimen \emph{H6} $w_v = 0{,}3$,
$w_t = 0{,}7$) sobre las mismas $101$ imágenes de test y $54$
clases.

Adicionalmente, se midió el \emph{eje idioma puro} sobre el propio
\emph{checkpoint} v4 de Dermapixel R0: el mismo modelo evaluado con
los \emph{prompts} castellanos originales frente a los mismos
\emph{prompts} traducidos al inglés, aislando el efecto del idioma
de consulta sobre una alineación cross-modal fija.

\section{Resultados}
\label{sec:anexoi-resultados}

La tabla~\ref{tab:ablation_panderm_en} reúne las seis
configuraciones evaluadas sobre el mismo conjunto de test.

\begin{table}[H]
\centering
\footnotesize
\setlength{\tabcolsep}{4pt}
\renewcommand{\arraystretch}{0.95}
\caption{Comparativa cuantitativa de Dermapixel R0 v4 (\emph{headline})
frente a la ablación con modelo fundacional inglés DermLIP v2 más
traducción ES--EN. Test \emph{holdout} DermapixelAI~1.0,
$N = 101$, $54$ clases L3 representadas; régimen \emph{H6 hybrid}
$v30\_r70$ salvo donde se indica.}
\label{tab:ablation_panderm_en}
\begin{tabular}{llrrr}
\hline
Sistema & Régimen & acc@1 L3 & BAcc L3 & acc@5 L3 \\
\hline
Dermapixel R0 v4    & H6 hybrid (castellano)            & $0{,}2475$          & $0{,}2685$ & $0{,}4455$ \\
Dermapixel R0 v5    & H6 hybrid (castellano)            & $0{,}2079$          & $0{,}2185$ & $0{,}4158$ \\
Dermapixel R0 v4    & H6 hybrid (inglés)              & $0{,}2277$          & $0{,}2556$ & $0{,}4554$ \\
DermLIP v2          & H6 hybrid (inglés)              & $0{,}2871$          & $0{,}3667$ & $0{,}6436$ \\
\textbf{DermLIP v2} & \textbf{templates+rich (inglés)} & $\mathbf{0{,}3861}$ & $0{,}4099$ & $0{,}6931$ \\
DermLIP v2          & visual-only                     & $0{,}2178$          & $0{,}2278$ & $0{,}4851$ \\
\hline
\end{tabular}
\end{table}

\section{Interpretación}
\label{sec:anexoi-interpretacion}

DermLIP v2 inglés supera a Dermapixel R0 v4 en el régimen
comparable \emph{H6} ($0{,}2871$ frente a $0{,}2475$,
$+3{,}96$\,pp) y lo supera con mayor margen en su mejor régimen
de plantillas más texto rico ($0{,}3861$, $+13{,}86$\,pp),
confirmando la hipótesis de partida de que el cuello de botella de
Dermapixel R0 residía en la \emph{escala} de la alineación
cross-modal y no en la formulación textual. Un modelo fundacional
preentrenado sobre $\approx 403\,000$ pares supera a un CLIP
castellano dedicado entrenado sobre $854$ pares, incluso pagando el
ruido de la traducción de etiquetas.

El eje idioma matiza esta lectura: sobre el mismo \emph{checkpoint}
v4, la consulta en castellano supera a la consulta en inglés
($0{,}2475$ frente a $0{,}2277$ en \emph{H6}, $+1{,}98$\,pp;
y $0{,}1287$ frente a $0{,}0891$ en \emph{rich-only},
$+3{,}96$\,pp), mientras que el régimen visual-only permanece
idéntico ($0{,}1683$). Dermapixel R0 sí captura un componente
específico del registro clínico castellano, si bien de segundo orden
frente a la ventaja de escala del modelo fundacional inglés.

El \emph{trade-off} entre los dos enfoques es, por tanto,
explícito:

\begin{itemize}
  \item \textbf{DermLIP inglés más traducción}: rendimiento bruto
        superior, cero entrenamiento, pero dependencia de
        traducción de etiquetas y de un dataset de preentrenamiento
        externo (Derm1M~\cite{derm1m}), con la consiguiente
        pérdida potencial de matices clínicos castellanos no
        traducibles uno a uno.
  \item \textbf{Dermapixel R0 (rama contrastiva)}: entrenamiento dedicado sobre $854$
        pares reales castellanos, preservación de la
        \emph{fingerprint} clínica del archivo de la
        Dra.~R.~Taberner, inferencia sin traducción en
        \emph{runtime} y operación \emph{on-prem} compatible con la
        ruta de extensión al banco hospitalario HUSLL
        (sección~\ref{sec:linea-active-learning}).
\end{itemize}

La conclusión es que la rama contrastiva de Dermapixel R0 no se
justifica como sistema de
mayor rendimiento bruto en \emph{zero-shot} L3, sino por sus
propiedades operativas de privacidad, especificidad lingüística y
autonomía de despliegue. La superioridad de DermLIP refuerza,
además, el diseño de la línea de continuación documentada en la
sección~\ref{sec:linea-active-learning}: la palanca pendiente
para Dermapixel R0 es la anotación experta dirigida
\emph{in-domain} sobre las clases L3 con peor cobertura, no la
escala genérica.

\section{Limitaciones de la ablación}
\label{sec:anexoi-limitaciones}

La comparación no está equiparada en datos: DermLIP v2 se
preentrena sobre $\approx 403\,000$ pares frente a los $854$ pares
castellanos de Dermapixel R0, y los \emph{backbones} visuales
difieren (PanDerm-Base contrastivo frente a PanDerm-Large
congelado con LoRA). El resultado debe leerse como ``modelo
fundacional inglés grande más traducción'' frente a ``CLIP
castellano dedicado pequeño'', no como un contraste de igual
presupuesto de datos.

La carga de DermLIP v2 requiere el \emph{fork} de
\texttt{open\_clip} del proyecto Derm1M: la implementación estándar
instancia un \emph{Vision Transformer} genérico e ignora
silenciosamente la torre visual PanDerm-Base ($152$ claves
ausentes, codificador visual aleatorio), un modo de fallo
verificado y evitado en esta ablación empleando el mismo
\emph{loader} que el módulo de producción
(Anexo~\ref{cap:anexoh}).

La verificación de solapamiento con Derm1M se realizó a nivel de
fuente: las fuentes constituyentes de Derm1M (vídeo, foros,
literatura biomédica y libros de texto en inglés y chino) no
incluyen ningún blog dermatológico español, y no se hallaron
coincidencias para \emph{dermapixel} ni \emph{taberner} en las
$417\,257$ filas del \emph{manifest}. El archivo Dermapixel es
privado y no figura como fuente, lo que sustenta razonablemente la
asunción de ausencia de \emph{leakage}; no se ejecutó una
comparación perceptual exhaustiva imagen a imagen. La traducción
ES--EN introduce, finalmente, un ruido no cuantificado en las
etiquetas.

\section{Detalles de reproducibilidad de la rama contrastiva}
\label{sec:anexoi-repro-contrastiva}

Este apartado recoge los detalles de procedencia y reproducibilidad
de la rama contrastiva de Dermapixel R0
(sección~\ref{sec:spanderm-clip}), separados del cuerpo por no
afectar a las conclusiones científicas.

\paragraph{Generación de \emph{captions} sintéticas.}
Las ocho perspectivas clínicas por imagen se generan mediante un
modelo de lenguaje grande a partir del texto original de cada caso
del blog y de sus metadatos. El total alcanza $7\,680$
\emph{captions} sobre $960$ imágenes (fase v1--v4); $102$ imágenes
fallan la generación automática y conservan únicamente el texto
base. El coste real de generación es de $32{,}54$ USD ($0{,}030$ USD
por imagen completa) y la longitud media es de $87{,}5$ palabras por
\emph{caption}. Combinadas con las nueve variantes textuales base
(plantillas L1/L2/L3, jerárquicas \texttt{L1L2L3}, título de blog,
tres \emph{chunks} de \texttt{case\_text} y \texttt{diagnosis\_raw}),
el dataset de entrenamiento alcanza del orden de $22\,000$ pares
(imagen, texto), muestreados por hash determinista
\texttt{(seed, epoch, image\_id)}.

\paragraph{Tiempos de entrenamiento.}
Sobre NVIDIA DGX Spark GB10: $\sim 18$\,min (v1), $\sim 3$\,min
(v2), $\sim 5$\,min (v3), $\sim 9$\,min (v4) y $\sim 54$\,min (v5).

\paragraph{Lección metodológica: longitud de las \emph{captions}.}
La generación de \emph{captions} sintéticas del dataset externo de
v5 (coste $42{,}33$ USD) reveló la importancia de calibrar la
longitud objetivo respecto al umbral del verificador automático
posterior. Un primer intento con rango $30$--$40$ palabras produjo
una tasa de rechazo del $54\%$ ---las \emph{captions} se agrupaban
en el suelo del rango y el verificador descartaba las más cortas---;
recalibrar el objetivo a $40$--$55$ palabras redujo el rechazo a
cero. La regla operativa derivada es dimensionar la longitud
objetivo con margen explícito sobre el umbral mínimo aceptable, no
sobre la media deseada.
